% !TEX program = pdflatex
\documentclass[12pt]{article}
\usepackage{geometry}
\geometry{left=2cm, right=2cm, top=2cm, bottom=2cm}
\usepackage{newtxtext}
\usepackage{hyperref}
\usepackage{setspace}
\onehalfspacing
\usepackage{parskip}

\pagestyle{empty}

\begin{document}

Dear Editor,

We are pleased to submit our manuscript, ``\textbf{\textit{Cooperation in competition? A biform game model for waste concrete recycling considering the cap-and-trade mechanism}},'' for consideration for publication in the \textit{IEEE Transactions on Engineering Management}.

This study develops a two-stage biform game model for waste concrete recycling that simultaneously integrates R\&D investment cooperation and non-cooperative pricing competition under the influence of a cap-and-trade mechanism. While existing research typically treats competition and cooperation as isolated stages, leaving it unclear whether enterprises should cooperate within a competitive environment, our model addresses this gap. By comparing our framework with a pure non-cooperative game model, we find that the biform game incentivizes a higher waste concrete conversion rate and expands total collection quantities. Crucially, cooperation achieves a Pareto improvement, yielding higher individual profits for both the manufacturer and the recycler while providing an effective buffering effect against the cost volatility raised by the cap-and-trade mechanism.

Our work aligns with the core scope of the \textit{IEEE Transactions on Engineering Management}, which emphasizes studies carried out within organizations to support decision-making and policy formation for RD\&E. Specifically, we investigate whether and how concrete manufacturers and recyclers should engage in backend R\&D investment cooperation within a front-end price-competitive market. Furthermore, we evaluate the operational impacts of environmental policies like the cap-and-trade mechanism on the participants' returns and economic decisions. The focuses on corporate technological R\&D choices and policy responses match the journal's publication criteria.

In terms of theoretical contributions, this study first constructs a two-stage biform game framework for waste concrete recycling, expanding the application boundaries of biform game theory. The framework effectively captures the intertwined dynamics of front-end pricing competition and back-end R\&D investment cooperation, overcoming the limitations of traditional single-stage non-cooperative games and offering a concrete analytical paradigm for understanding horizontal coopetition. Second, this research innovatively integrates carbon emission reduction and the cap-and-trade mechanism into the economic decisions of waste concrete recycling management. By coupling carbon market dynamics with cooperative investment strategies, we theoretically demonstrate how cooperation can transform carbon compliance from a static constraint into a driver for R\&D investment.

Practically, the study provides actionable guidance for industry practitioners and policymakers. First, concrete manufacturers and recyclers should actively seek R\&D cooperation, as pure price competition in the collection market fails to maximize mutual economic interests. Second, enterprise managers should evaluate the operational impacts of environmental policies like the cap-and-trade mechanism. To effectively buffer against carbon market volatility and high compliance costs, firms are advised to implement KPI-linked technology cost-sharing contracts and establish internal carbon risk-sharing mechanisms, transforming regulatory constraints into cooperative innovation drivers. Finally, government policymakers should simultaneously weigh enterprise economic viability and long-term environmental development when setting carbon emission caps and trading prices. Furthermore, the government can implement targeted structural subsidies like the logistical subsidy and incentivize R\&D investments cooperation, fostering a sustainable and cooperative ecosystem between concrete manufacturers and recyclers in the waste concrete industry.

We believe this work is timely and relevant to IEEE Transactions on Engineering Management's readership. The manuscript is not under review elsewhere, has been approved by all co-authors, and presents no conflicts of interest.

Thank you very much for considering our work. We look forward to your feedback.

Yours sincerely,

Zilun Wang\\
Suzhou university of science and technology\\
Suzhou, China\\
Email: zylenw97@usts.edu.cn

\end{document}